\chapter*{Abstract}
\addcontentsline{toc}{chapter}{Abstract}
\markboth{Abstract}{Abstract}

\esibot{} is a differential-drive autonomous mobile robot designed and built by the team
as part of the IoT integrative project at the École Nationale Supérieure d'Informatique
(ESI Algiers). The initial specification was inspired by the TurtleBot3 Burger
(ROBOTIS): an ESP32-CAM was meant to directly drive the motors, the servo motor and a
rotating ultrasonic radar (HC-SR04 + SG90), while the Raspberry\,Pi\,4 hosted the \ros{}
software stack (\gls{slam}, Nav2, web dashboard).

Over the course of development, the hardware architecture actually implemented evolved
significantly from this specification. Control of the two DC motors, encoder readout and
velocity control (\gls{pid}) are now handled \emph{directly} by the \gls{gpio} pins of the
Raspberry\,Pi\,4, assisted by an MPU6050 inertial measurement unit fused through an
extended Kalman filter (\texttt{robot\_localization}). The ESP32-CAM is retained solely
as a \gls{wifi} streaming camera (\gls{mjpeg}), no longer participating in actuator
control. The rotating ultrasonic radar remains available as a slow fallback path, but
current mapping relies on a Kinect (v1, Xbox\,360 era) depth sensor, driven by
\texttt{kinect\_ros2} and converted to a laser scan via
\texttt{depthimage\_to\_laserscan}. The \texttt{working\_build} branch also adds the
\texttt{depth\_cleanup} library to prepare a cleaned depth stream, and references
\texttt{m-explore-ros2} for autonomous exploration.

This report faithfully documents this as-built architecture: mechanical design
(differential chassis described by the \gls{urdf} model), electronic wiring inferred
from code parameters, low-level control and motor control loops, the perception
pipeline (\gls{imu}, Kinect), \gls{slam} mapping, autonomous \ros{} Nav2 navigation, a
telemetry web dashboard, and Xbox controller teleoperation. Measured results (or results
derived analytically from the parameters declared in the code, where field measurements
were still in progress) are presented, followed by a critical analysis of the
discrepancies between the original design and the actual implementation, the identified
limitations, and an outlook for future improvements.

\vspace{6mm}
\noindent\textbf{Keywords:} \ros{} Humble, differential-drive robot, \gls{slam},
slam\_toolbox, Nav2, Kinect, Raspberry\,Pi\,4, ESP32-CAM, MPU6050, extended Kalman filter,
autonomous navigation, IoT robotics, TurtleBot3.
